% Source: https://artofproblemsolving.com/community/c3h2493539_2018_mathcounts_states_23
% Stable identity: aops-topic-2493539-q01
Carl can always open five empty boxes among the nine that Bryan did not take, whatever is in Bryan's box, so his action does not depend on whether Bryan's box holds the coin.
\begin{itemize}
\item With probability $\dfrac1{10}$ Bryan's first box has the coin; then switching certainly loses.
\item With probability $\dfrac9{10}$ his first box is empty. Then the coin is in one of the $9$ other boxes, and since Carl opened $5$ empty ones, it is in one of the $10-1-5=4$ remaining closed boxes. Bryan picks one of these $4$ at random and wins with probability $\dfrac14$.
\end{itemize}
So switching wins with probability
\[\dfrac9{10}\cdot\dfrac14=\boxed{\dfrac{9}{40}}.\]
\emph{Why the first box stays at $\dfrac1{10}$.} Opening empty boxes that Carl was always able to open gives no information about Bryan's own box, so it still holds the coin with probability $\dfrac1{10}=\dfrac4{40}$. The remaining probability $\dfrac9{10}$ now sits on only $4$ closed boxes, $\dfrac{9}{40}$ each, which is more than $\dfrac4{40}$: switching is better, but the chance of picking the right box is still less than $\dfrac14$, because with probability $\dfrac1{10}$ the coin was in the returned box.
